\section*{Capítulo \ref{ch:b2:polymer-synthesis} --- Polímeros: síntese, estrutura e propriedades}

\begin{solution}{exo:b2:polymer-synthesis:1}
Polieteno \ce{-[CH2-CH2]-}; polipropeno \ce{-[CH2-CH(CH3)]-}; poli(cloreto de vinila)
\ce{-[CH2-CHCl]-}; poliestireno \ce{-[CH2-CH(C6H5)]-}; náilon-6,6 \ce{-[NH(CH2)6NH-CO(CH2)4CO]-}.
\end{solution}

\begin{solution}{exo:b2:polymer-synthesis:2}
PET: crescimento em etapas (poliéster, água liberada). Poliestireno: crescimento em cadeia. Náilon-6,6:
crescimento em etapas. Poli(metacrilato de metila): crescimento em cadeia (um \omterm{def:b2:polymer-synthesis:polymer}{monômero} com \ce{C=C}).
Poli(ácido lático) a partir do hidroxiácido: crescimento em etapas, um \omterm{def:b2:polymer-synthesis:polymer}{monômero} \ce{A-B}.
\end{solution}

\begin{solution}{exo:b2:polymer-synthesis:3}
$\bar M_n = 1000 \times \qty{104.15}{g/mol} \approx \qty{104}{kg/mol}$ (estireno \ce{C8H8},
\qty{104.152}{g/mol}).
% ledger: aw:C, aw:H
\end{solution}

\begin{solution}{exo:b2:polymer-synthesis:4}
\omterm{def:b2:polymer-synthesis:tacticity}{Sindiotática}. Sim: uma cadeia regular pode se empacotar em cristalitos.
\end{solution}

\begin{solution}{exo:b2:polymer-synthesis:5}
Frações mássicas 0.5 e 0.5. $1/\bar M_n = 0.5/10 + 0.5/100 = 0.055$, $\bar M_n \approx
\qty{18.2}{kg/mol}$; $\bar M_w = 0.5 \times 10 + 0.5 \times 100 = \qty{55}{kg/mol}$; $\text{\DH} \approx
3.0$.
\end{solution}

\begin{solution}{exo:b2:polymer-synthesis:6}
$1/(1-0.98) = 50$; $1/(1-0.995) = 200$.
\end{solution}

\begin{solution}{exo:b2:polymer-synthesis:7}
$r = 1/1.01$; at $p = 1$, $\bar X_n = (1+r)/(1-r) = (1.01 + 1)/(1.01 - 1) = 201$.
\end{solution}

\begin{solution}{exo:b2:polymer-synthesis:8}
$R_p \propto [\mathrm I]^{1/2}$: multiplicada por $\sqrt2 \approx 1.41$. $\nu \propto [\mathrm I]^{-1/2}$:
dividido por $\sqrt 2$, cerca de 0.71 vez seu valor anterior.
\end{solution}

\begin{solution}{exo:b2:polymer-synthesis:9}
O PVC puro tem sua transição vítrea acima da temperatura ambiente: um vidro rígido. As pequenas
moléculas de plastificante se instalam entre as cadeias, as afastam e permitem que os segmentos se movam a
temperatura mais baixa: a $T_g$ cai abaixo da temperatura ambiente, e o material é borrachoso e flexível
no uso.
\end{solution}

\begin{solution}{exo:b2:polymer-synthesis:10}
$\dfrac{\mathrm d}{\mathrm dx}\bigl(x p^{\,x-1}\bigr) = p^{\,x-1}(1 + x\ln p)$, nula para $x = -1/\ln p$
(um máximo: o sinal passa de $+$ a $-$). Para $p$ próximo de 1, $\ln p \approx -(1-p)$, logo $x \approx
1/(1-p) = \bar X_n$.
\end{solution}

\begin{solution}{exo:b2:polymer-synthesis:11}
Fase 1, \omterm{def:b2:acyl-substitution:transesterification}{transesterificação}: \ce{C6H4(COOCH3)2 + 2HOCH2CH2OH -> C6H4(COOCH2CH2OH)2 + 2CH3OH}, metanol
removido por destilação. Fase 2: cada molécula de diéster leva duas extremidades \ce{CH2CH2OH}; uma
extremidade ataca um éster de outra molécula e libera uma molécula de etano-1,2-diol, que é bombeada
para fora. O \omterm{def:b2:polymer-synthesis:polymer}{monômero} leva os dois tipos de grupo na proporção certa, como um \omterm{def:b2:polymer-synthesis:polymer}{monômero} \ce{A-B}: o
equilíbrio $r = 1$ está embutido, e remover o diol leva $p$ para 1.
\end{solution}

\begin{solution}{exo:b2:polymer-synthesis:12}
Alimentação equimolar: razão $= (2.0 + 1)/(1 + 0.5) = 2$: duas vezes mais unidades do \omterm{def:b2:polymer-synthesis:polymer}{monômero} 1. Com $r_1 = r_2 =
0$, razão $= [\mathrm M_1][\mathrm M_2]/([\mathrm M_2][\mathrm M_1]) = 1$ qualquer que seja a alimentação: cada
radical adiciona apenas o outro \omterm{def:b2:polymer-synthesis:polymer}{monômero}, um \omterm{def:b2:polymer-synthesis:copolymer}{copolímero} alternado.
\end{solution}

\begin{solution}{pb:b2:polymer-synthesis:1}
\textbf{1.} \ce{H2N(CH2)6NH2} e \ce{HOOC(CH2)4COOH}.
\textbf{2.} Uma água por ligação:
\[\mathrm{{-}COOH + H_2N{-} \longrightarrow {-}CO{-}NH{-} + H_2O}.\]
\textbf{3.} O sal contém exatamente uma diamina por diácido; pesar separadamente dois líquidos ou
sólidos nunca atingiria a exatidão que a equação de Carothers exige.
\textbf{4.} \ce{-[NH(CH2)6NH-CO(CH2)4CO]-}, \ce{C12H22N2O2}, \qty{226.32}{g/mol}.
\textbf{5.} A \omterm{def:b2:polymer-synthesis:polymer}{unidade de repetição} contém duas unidades provenientes de \omterm{def:b2:polymer-synthesis:polymer}{monômeros}: $M_0 = 226.32/2 = \qty{113.16}{g/mol}$.
\textbf{6.} Os carbonos da diamina (6) e do diácido (6).
\textbf{7.} $\bar X_n = 50$, $\bar M_n = 50 \times 113.16 \approx \qty{5.66}{kg/mol}$.
\textbf{8.} $\bar X_n = 100$, $\bar M_n \approx \qty{11.3}{kg/mol}$.
\textbf{9.} Com 99~\% dos grupos tendo reagido, a cadeia média tem só 100 unidades, curta demais para
uma fibra resistente; cada passo a mais rumo à \omterm{def:b2:continuous-reactors:conversion}{conversão} completa dobra ou triplica o comprimento da
cadeia.
\textbf{10.} $\bar X_n = 10\,000/113.16 = 88.4$; $p = 1 - 1/88.4 \approx 0.989$.
\textbf{11.} Removendo a água: o material fundido é aquecido bem acima de seu ponto de fusão, por fim sob
pressão reduzida, de modo que o equilíbrio continue se deslocando para a amida.
\textbf{12.} Uma impureza monofuncional bloqueia extremidades de cadeia, e toda impureza perturba o
equilíbrio 1\,:\,1.
\textbf{13.} $r = 1/1.01 \approx 0.990$.
\textbf{14.} $\bar X_n = (1+r)/(1-r) = 201$; $\bar M_n \approx 201 \times 113.16 \approx
\qty{22.7}{kg/mol}$.
\textbf{15.} $\bar X_n = 1.9901/(1.9901 - 2 \times 0.9901 \times 0.99) \approx 1.9901/0.0297 \approx
67$.
\textbf{16.} Grupos amina: quando os grupos ácidos se esgotam, toda cadeia termina em \ce{NH2}.
\textbf{17.} Ele transforma uma extremidade amina em uma amida que não pode mais reagir: ele bloqueia
as cadeias, limitando $\bar M_n$ e impedindo um crescimento posterior quando o \omterm{def:b2:polymer-synthesis:polymer}{polímero} é refundido.
\textbf{18.} Para manter uma viscosidade do fundido reprodutível para a fiação e a moldagem, que de
outro modo derivaria à medida que as cadeias continuam reagindo no fundido.
\textbf{19.} $\text{\DH} = 1 + p = 1.99$.
\textbf{20.} $\bar M_w = 1.99 \times 11.3 \approx \qty{22.5}{kg/mol}$.
\textbf{21.} Fração em número $n_1 = 1 - p = 0.01$, 1~\% das moléculas; fração mássica $w_1 = (1-p)^2 =
10^{-4}$, 0.01~\% da massa.
\textbf{22.} Perto de $x = -1/\ln 0.99 \approx 99.5$, isto é, perto de $\bar X_n = 100$ (exercício 10).
\textbf{23.} $1000/226.32 = \qty{4.42}{mol}$ de \omterm{def:b2:polymer-synthesis:polymer}{unidades de repetição}, duas ligações amida cada: \qty{8.84}{mol} de
água, $8.84 \times 18.015 \approx \qty{159}{g}$.
\textbf{24.} Cadeias curtas demais dão uma fibra fraca; cadeias longas demais dão um fundido viscoso
demais para ser empurrado pelos furos finos da fieira.
\textbf{25.} $\bar X_n = 20\,000/113.16 = 176.7$ e $p = 1 - 1/176.7 = \boldsymbol{\approx 0.994}$.
\end{solution}
